\documentclass[11pt,a4paper]{report}

% Essential packages
\usepackage[utf8]{inputenc}
\usepackage[T1]{fontenc}
\usepackage{lmodern}
\usepackage[english]{babel}
\usepackage{graphicx}
\usepackage{xcolor}
\usepackage{geometry}
\usepackage{fancyhdr}
\usepackage{titlesec}
\usepackage{tcolorbox}
\usepackage{enumitem}
\usepackage{booktabs}
\usepackage{longtable}
\usepackage{multirow}
\usepackage{array}
\usepackage{listings}
\usepackage{hyperref}
\usepackage{tikz}
\usepackage{pgfgantt}
\usepackage[most]{tcolorbox}

% Page geometry
\geometry{margin=2.5cm}

% Colors
\definecolor{primary}{RGB}{0,82,155}
\definecolor{secondary}{RGB}{100,100,100}
\definecolor{highlight}{RGB}{242,93,34}
\definecolor{light}{RGB}{240,240,240}
\definecolor{success}{RGB}{40,167,69}
\definecolor{warning}{RGB}{255,193,7}

% Header and footer
\pagestyle{fancy}
\fancyhf{}
\fancyhead[L]{\textcolor{primary}{\leftmark}}
\fancyhead[R]{\textcolor{secondary}{Final Project Review}}
\fancyfoot[C]{\thepage}
\renewcommand{\headrulewidth}{0.4pt}
\renewcommand{\footrulewidth}{0.4pt}

% Title formatting
\titleformat{\chapter}[display]
{\normalfont\huge\bfseries\color{primary}}
{\chaptertitle}{20pt}{\Huge}
\titlespacing*{\chapter}{0pt}{-50pt}{40pt}

\titleformat{\section}
{\color{primary}\normalfont\Large\bfseries}
{\thesection}{1em}{}

\titleformat{\subsection}
{\color{secondary}\normalfont\large\bfseries}
{\thesubsection}{1em}{}

% Custom environments
\newenvironment{status}[2][]{
    \begin{tcolorbox}[
        enhanced,
        attach boxed title to top center={yshift=-3mm},
        colback=white,
        colframe=#1!30,
        colbacktitle=#1!30,
        title=Status: #2,
        fonttitle=\bfseries
    ]
}{
    \end{tcolorbox}
}

% Code styling
\lstset{
    basicstyle=\ttfamily\small,
    backgroundcolor=\color{light},
    frame=single,
    numbers=left,
    numberstyle=\tiny\color{gray},
    keywordstyle=\color{blue},
    commentstyle=\color{secondary},
    stringstyle=\color{highlight},
    breaklines=true,
    breakatwhitespace=true,
    showspaces=false,
    showstringspaces=false,
    showtabs=false
}

% Progress bar command
\newcommand{\progressbar}[2][primary]{
    \begin{tikzpicture}
        \fill[#1!20] (0,0) rectangle (5,0.4);
        \fill[#1] (0,0) rectangle (#2*5/100,0.4);
        \node[right] at (5.2,0.2) {\textbf{#2\%}};
    \end{tikzpicture}
}

\begin{document}

% Title page
\begin{titlepage}
\begin{center}
    \vspace*{2cm}
    \begin{tcolorbox}[
        enhanced,
        width=\textwidth,
        colback=white,
        colframe=primary,
        arc=0mm,
        boxrule=2pt
    ]
    \begin{center}
        \vspace{1cm}
        {\Huge\textcolor{primary}{\textbf{JPF Stretch Hub}}\par}
        \vspace{0.5cm}
        {\Large\textcolor{secondary}{Final Project Review Report}\par}
        \vspace{1cm}
    \end{center}
    \end{tcolorbox}
    
    \vspace{2cm}
    {\large\textbf{Prepared by:}\par}
    \vspace{0.5cm}
    {\Large Meshack Mogire\par}
    
    \vspace{1cm}
    {\large\textbf{Date:}\par}
    \vspace{0.5cm}
    {\Large March 5, 2025\par}
    
    \vspace{2cm}
    \begin{tcolorbox}[
        enhanced,
        width=0.6\textwidth,
        colback=success!10,
        colframe=success,
        arc=0mm
    ]
    \begin{center}
        \textbf{Overall Project Status:} 90\% Complete\\
        \vspace{0.3cm}
        \progressbar[success]{90}
    \end{center}
    \end{tcolorbox}
    
    \vfill
\end{center}
\end{titlepage}

% Table of contents
\tableofcontents
\clearpage

\chapter{Project Overview}

\section{Application Description}

JPF Stretch Hub is a comprehensive web application for bookings, payments, and user management at a stretching and fitness studio. The platform enables clients to book sessions with stretch instructors, manage their packages, and make payments. Staff can manage schedules, services, and client information through a unified interface.

\section{Technologies Used}

\subsection{Backend Stack}
\begin{tcolorbox}[
    enhanced,
    colback=light,
    colframe=primary,
    arc=0mm
]
\begin{itemize}
    \item Django 5.0 with Django REST Framework
    \item PostgreSQL for database management
    \item Celery with Redis for asynchronous task processing
    \item Django Allauth with MFA support for authentication
\end{itemize}
\end{tcolorbox}

\subsection{Frontend Technologies}
\begin{tcolorbox}[
    enhanced,
    colback=light,
    colframe=primary,
    arc=0mm
]
\begin{itemize}
    \item HTML/CSS with Tailwind CSS for styling
    \item AlpineJS for reactive components and interactivity
    \item Responsive design principles for mobile compatibility
\end{itemize}
\end{tcolorbox}

\subsection{DevOps \& Infrastructure}
\begin{tcolorbox}[
    enhanced,
    colback=light,
    colframe=primary,
    arc=0mm
]
\begin{itemize}
    \item Docker for containerization
    \item Redis for caching and message brokering
\end{itemize}
\end{tcolorbox}

\section{Modular Structure}

The application is organized into distinct, interconnected modules:

\begin{longtable}{>{\raggedright}p{0.25\textwidth}>{\raggedright\arraybackslash}p{0.65\textwidth}}
\toprule
\textbf{Module} & \textbf{Purpose} \\
\midrule
\endhead
users & Custom user model with profile management \\
bookings & Booking system for scheduling sessions \\
services & Service offerings and packages management \\
locations & Physical locations management \\
stretch\_instructors & Instructor profile and availability management \\
clients & Client management and history tracking \\
schedules & Scheduling system for sessions \\
packages & Service package management and tracking \\
\bottomrule
\end{longtable}

\chapter{Areas to Verify and Finalize}

\section{Validation and Error Handling}
\begin{status}[warning]{70\% Complete}
The validation system currently exists across models, forms, and views but lacks centralization and comprehensiveness. A more robust approach is needed to ensure all user inputs are properly validated and appropriate error messages are displayed.

\subsection*{Current Implementation}
\begin{itemize}
    \item Basic field validation in Django models
    \item Form validation for required fields
    \item Some business logic validation spread across views
\end{itemize}

\subsection*{Required Improvements}
\begin{itemize}
    \item Centralized validation strategy
    \item Comprehensive field validation
    \item Consistent error handling and display
\end{itemize}
\end{status}

% ... Continue with the rest of the sections ...

\end{document}
